\documentclass[exercice]{thedocument}%
\usepackage{tabularray}%
\startdatakeys
\source{}
\cycle{}
\competencesDuSocle{}
\connaissancesRequises{}
\competencesTravaillees{}
\enddatakeys
\begin{document}%
Les tableaux suivants traduisent une situation de proportionnalité.
\vskip5pt\par\noindent
Les compléter en utilisant la méthode la plus adaptée pour chacun d'entre
eux.\vskip10pt\noindent
\hbox{\hbox to.5\linewidth{{\bfseries a)\hskip10pt}%
\begin{tblr}{colspec={*2{Q[c,m]}},hlines,vlines}
    \<n1;&\<n3;\\
    \<n2;&
\end{tblr}
}\hbox{{\bfseries b)\hskip10pt}%
\begin{tblr}{colspec={*3{Q[c,m]}},hlines,vlines}
    \<n4;&\<n6;&\\
    \<n5;&&\<n7;
\end{tblr}
}}\vskip10pt\noindent
\hbox{\hbox to.5\linewidth{{\bfseries c)\hskip10pt}%
\begin{tblr}{colspec={*3{Q[c,m]}},hlines,vlines}
    \<n8;&\<n10;&\<n12;\\
    \<n9;&\<n11;&
\end{tblr}
}\hbox{{\bfseries d)\hskip10pt}%
\begin{tblr}{colspec={*5{Q[c,m]}},hlines,vlines}
    \<n13;&&&&\\
    \<n14;&\<n15;&\<n16;&\<n17;&\<n18;
\end{tblr}
}}
\end{document}
